\documentclass[12pt]{article}
\usepackage[utf8]{inputenc}
\usepackage[T1]{fontenc}
\usepackage{hyperref}
\title{Ketengban Language Resources for K-12 Curriculum Development}
\author{Compiled by Claude Sonnet 5.5 (claude-sonnet-5-5) for Andr\'as Kornai}
\date{}
\begin{document}
\maketitle

\section{Introduction}
Ketengban (also Kupel, Oktengban; ISO 639-3 xte) is a Trans-New Guinea language of the Eastern Mek branch. It has about 10,000 speakers in Pegunungan Bintang, Highland Papua, Indonesia, near the Papua New Guinea border. Listed varieties are Okbap, Omban, Bime, Onya and Sirkai. Very little published educational material was found in the searches. This dossier lists only what the searches returned.

\section{Primary Educational Resources}
No primers, textbooks or school programmes for Ketengban turned up in the searches, and the language is not known to be taught in schools.

\section{Linguistic and Reference Materials}
\begin{itemize}
\item Wikipedia, ``Ketengban language'': a general overview. \url{https://en.wikipedia.org/wiki/Ketengban_language}
\item Omniglot, Ketengban: notes on the Latin-based writing and the language. \url{https://omniglot.com/writing/ketengban.htm}
\item New Guinea World (University of Zurich), Ketengban entry, with classification and bibliographic leads. \url{https://newguineaworld.linguistik.uzh.ch/families/trans-new-guinea/central-west-new-guinea/momuna-mek/mek/east-mek/ketengban}
\item Glosbe Ketengban--English: a community-built word list. It is not a scholarly dictionary and needs checking by speakers. \url{https://glosbe.com/xte/en}
\end{itemize}
No dedicated Ketengban grammar was found.

\section{Religious and Community Resources}
\begin{itemize}
\item Scripture Earth lists Ketengban scripture texts with audio and the JESUS Film. Sources disagree on whether the translation covers the whole Bible or the New Testament only, so check the listing. \url{https://www.scriptureearth.org/00eng.php?iso=xte}
\item Joshua Project, Kupel/Ketengban in Indonesia: demographic information. \url{https://joshuaproject.net/people_groups/12866/ID}
\end{itemize}

\section{Digital and Audio Resources}
The Scripture Earth page above provides audio and app access to the scripture text. No other corpora or software were found.

\section{Curriculum Development Context}
Material in the language is very limited: scripture, a community word list and a few overview pages. A curriculum would need to begin with community-led documentation: orthography confirmation, story collection and local teacher training. Indonesian is the likely language of instruction in the interim. The Bible text could serve as an initial reading corpus, with caution about register and content.
\end{document}